\documentclass{article}
\usepackage{amsmath,amssymb,amsthm}
\usepackage{algorithm2e}
\usepackage{hyperref}
\usepackage{enumitem}
\newtheorem{theorem}{Theorem}
\newtheorem{lemma}{Lemma}
\newtheorem{assumption}{Assumption}
\newcommand{\E}{\mathbb{E}}
\newcommand{\Var}{\mathrm{Var}}

\begin{document}

\section*{Cyclic Stochastic Gradient Descent (Cyclic SGD)}

\section*{Mathematical Formulation}
Let $f:\mathbb{R}^d\rightarrow\mathbb{R}$ be the (possibly non‑convex) objective and let 
$\{\,\xi_t\,\}_{t\ge 0}$ denote an i.i.d.\ sequence of random variables that generate stochastic
gradients 
\[
g_t \;:=\; \nabla f(w_t; \xi_t).
\]
The learning rate is \emph{cyclic} with period $M\in\mathbb{N}_{+}$ and amplitude 
$[\eta_{\min},\eta_{\max}]$, $0<\eta_{\min}\le \eta_{\max}$.  
We adopt the widely used \textbf{triangular} schedule
\[
\boxed{\;
\eta_t \;=\; \eta_{\min}+(\eta_{\max}-\eta_{\min})\,
\Bigl(1-\frac{\bigl|\, (t\bmod 2M)-M\,\bigr|}{M}\Bigr),\qquad t=0,1,2,\dots
\;}
\tag{1}
\]
which linearly increases from $\eta_{\min}$ to $\eta_{\max}$ in the first $M$ steps
and then linearly decreases back to $\eta_{\min}$ in the next $M$ steps, repeating
forever.

The update rule is
\[
\boxed{\; w_{t+1}= w_t-\eta_t\, g_t,\qquad t=0,1,2,\dots \;}
\tag{2}
\]

\section*{Pseudocode}
\begin{algorithm}[H]
\caption{Cyclic SGD}
\KwIn{Initial point $w_0\in\mathbb{R}^d$, period $M$, bounds $\eta_{\min},\eta_{\max}$,
maximum iterations $T$.}
\For{$t=0$ \KwTo $T-1$}{
    $\displaystyle \eta_t \leftarrow \eta_{\min}+(\eta_{\max}-\eta_{\min})
    \Bigl(1-\frac{|(t\bmod 2M)-M|}{M}\Bigr)$\;
    Sample $\xi_t$ and compute stochastic gradient $g_t=\nabla f(w_t;\xi_t)$\;
    $w_{t+1}\leftarrow w_t-\eta_t g_t$\;
}
\Return $w_T$
\end{algorithm}

\section*{Dry Run Example}
Consider the scalar quadratic $f(w)=(w-3)^2$, whose deterministic gradient is
$ \nabla f(w)=2(w-3)$.  Use a triangular schedule with $M=2$, $\eta_{\min}=0.05$,
$\eta_{\max}=0.15$, and start from $w_0=0$.

\begin{center}
\begin{tabular}{c|c|c|c|c}
$t$ & $\eta_t$ (from (1)) & $g_t=2(w_t-3)$ & $w_{t+1}=w_t-\eta_t g_t$ & $f(w_{t+1})$\\\hline
0 & $0.05$ & $2(0-3)=-6$ & $0-0.05(-6)=0.30$ & $(0.30-3)^2=7.29$\\
1 & $0.10$ & $2(0.30-3)=-5.40$ & $0.30-0.10(-5.40)=0.84$ & $(0.84-3)^2=4.66$\\
2 & $0.15$ & $2(0.84-3)=-4.32$ & $0.84-0.15(-4.32)=1.38$ & $(1.38-3)^2=2.62$
\end{tabular}
\end{center}
The iterates move toward the minimiser $w^\star=3$ while the step size follows the prescribed cycle.

\newpage
\section*{Assumptions}
\begin{assumption}[Smoothness]\label{as:smooth}
$f$ is $L$‑smooth, i.e.
\[
\|\nabla f(x)-\nabla f(y)\|\le L\|x-y\|\quad\forall x,y\in\mathbb{R}^d .
\]
\end{assumption}
\begin{assumption}[Unbiased Stochastic Gradient]\label{as:unbiased}
For every $t$, $\E[g_t\mid w_t]=\nabla f(w_t)$.
\end{assumption}
\begin{assumption}[Bounded Variance]\label{as:variance}
There exists $\sigma^2\ge0$ such that
\[
\E\big[\|g_t-\nabla f(w_t)\|^2\mid w_t\big]\le\sigma^2,\qquad\forall t .
\]
\end{assumption}
\begin{assumption}[Bounded Learning Rates]\label{as:lr}
The cyclic schedule satisfies
\[
0<\eta_{\min}\le \eta_t\le\eta_{\max}<\frac{1}{L},\qquad\forall t .
\]
\end{assumption}

\section*{Main Convergence Theorem}
\begin{theorem}[Average Gradient Norm for Cyclic SGD]\label{thm:main}
Let $\{w_t\}_{t\ge0}$ be generated by \eqref{1}–\eqref{2} under Assumptions
\ref{as:smooth}–\ref{as:lr}.  Define the weighted average step size
\[
\bar\eta_T \;:=\; \frac{1}{T}\sum_{t=0}^{T-1}\eta_t .
\]
Then for any $T\ge 1$,
\[
\frac{1}{T}\sum_{t=0}^{T-1}\E\big[\|\nabla f(w_t)\|^2\big]
\;\le\;
\frac{2\big(f(w_0)-f^\star\big)}{\bar\eta_T\,T}
\;+\;
\frac{L\sigma^2}{\bar\eta_T\,T}\sum_{t=0}^{T-1}\eta_t^{\,2},
\tag{3}
\]
where $f^\star:=\inf_{w}f(w)$.  In particular, because
$\eta_{\min}\le\eta_t\le\eta_{\max}$,
\[
\bar\eta_T\ge\eta_{\min},\qquad 
\frac{1}{T}\sum_{t=0}^{T-1}\eta_t^{\,2}\le\eta_{\max}^{2},
\]
so that
\[
\frac{1}{T}\sum_{t=0}^{T-1}\E\big[\|\nabla f(w_t)\|^2\big]
\;\le\;
\underbrace{\frac{2\big(f(w_0)-f^\star\big)}{\eta_{\min}T}}_{\displaystyle O\!\left(\frac1T\right)}
\;+\;
\underbrace{\frac{L\sigma^2\eta_{\max}^{2}}{\eta_{\min}T}}_{\displaystyle O\!\left(\frac1T\right)} .
\tag{4}
\]
Thus the method attains the optimal $O(1/T)$ rate for the squared gradient norm in the
non‑convex smooth setting.
\end{theorem}

\section*{Theoretical Properties}
\begin{enumerate}[label=\arabic*.]
\item \textbf{Stability.} Since $\eta_t<L^{-1}$ for all $t$, the standard descent
inequality for $L$‑smooth functions (Lemma~\ref{lem:descent}) holds uniformly, guaranteeing that the iterates remain bounded.
\item \textbf{Hyperparameter Sensitivity.}
The bound \eqref{3} depends only on the extremal learning rates
$\eta_{\min}$ and $\eta_{\max}$.  Increasing the amplitude
$\eta_{\max}-\eta_{\min}$ enlarges the second term in \eqref{4} but also
potentially accelerates early progress; the bound captures this trade‑off.
\item \textbf{Comparison to Baselines.}
For a constant step size $\eta_t\equiv\eta$, \eqref{3} reduces to the classical
SGD bound $\frac{2(f(w_0)-f^\star)}{\eta T}+L\sigma^2\eta/T$.
Cyclic SGD inherits the same $O(1/T)$ rate while allowing the practitioner
to explore larger step sizes intermittently, a phenomenon observed empirically.
\end{enumerate}

\section*{Discussion}
The theorem shows that, despite the non‑monotonic variation of the step size,
the average squared gradient norm decays at the same order as for standard SGD.
The proof hinges on the fact that every $\eta_t$ is bounded above by $1/L$,
so the usual smoothness inequality can be applied at each iteration.
The cyclic schedule therefore does not deteriorate the worst‑case theoretical
guarantee, while in practice it can help escape shallow local minima.

\newpage
\section*{Preliminaries (Definitions and Lemmas)}
\begin{lemma}[Descent Lemma for $L$‑smooth functions]\label{lem:descent}
If $f$ is $L$‑smooth, then for any $x,y\in\mathbb{R}^d$,
\[
f(y)\;\le\; f(x)+\langle\nabla f(x),y-x\rangle+\frac{L}{2}\|y-x\|^2 .
\]
\end{lemma}
\begin{proof}
Standard; follows from integrating the Lipschitz condition on the gradient.
\end{proof}

\section*{Main Proof (Step‑by‑Step Derivation)}
We prove Theorem~\ref{thm:main}.  Throughout the proof, expectations are taken
with respect to all randomness up to time $t$ unless otherwise noted.

\begin{enumerate}
\item[Step 1.] Apply Lemma~\ref{lem:descent} with $x=w_t$, $y=w_{t+1}=w_t-\eta_t g_t$:
\[
f(w_{t+1})\le f(w_t)-\eta_t\langle\nabla f(w_t),g_t\rangle
+\frac{L\eta_t^{\,2}}{2}\|g_t\|^2 .
\tag{5}
\]

\item[Step 2.] Expand the inner product using $g_t=\nabla f(w_t)+\delta_t$, where
$\delta_t:=g_t-\nabla f(w_t)$:
\[
\langle\nabla f(w_t),g_t\rangle
= \|\nabla f(w_t)\|^2+\langle\nabla f(w_t),\delta_t\rangle .
\tag{6}
\]

\item[Step 3.] Substitute \eqref{6} into \eqref{5}:
\[
f(w_{t+1})\le f(w_t)-\eta_t\|\nabla f(w_t)\|^2
-\eta_t\langle\nabla f(w_t),\delta_t\rangle
+\frac{L\eta_t^{\,2}}{2}\|g_t\|^2 .
\tag{7}
\]

\item[Step 4.] Expand $\|g_t\|^2=\|\nabla f(w_t)+\delta_t\|^2
= \|\nabla f(w_t)\|^2+2\langle\nabla f(w_t),\delta_t\rangle+\|\delta_t\|^2$.
Insert into \eqref{7}:
\[
\begin{aligned}
f(w_{t+1})\le &\, f(w_t)-\eta_t\|\nabla f(w_t)\|^2
-\eta_t\langle\nabla f(w_t),\delta_t\rangle \\
&\;+\frac{L\eta_t^{\,2}}{2}\Bigl(\|\nabla f(w_t)\|^2
+2\langle\nabla f(w_t),\delta_t\rangle+\|\delta_t\|^2\Bigr) .
\end{aligned}
\tag{8}
\]

\item[Step 5.] Collect the terms containing $\|\nabla f(w_t)\|^2$ and
$\langle\nabla f(w_t),\delta_t\rangle$:
\[
\begin{aligned}
f(w_{t+1})\le &\, f(w_t)
-\Bigl(\eta_t-\frac{L\eta_t^{\,2}}{2}\Bigr)\|\nabla f(w_t)\|^2 \\
&\;-\Bigl(\eta_t-\!L\eta_t^{\,2}\Bigr)\langle\nabla f(w_t),\delta_t\rangle
+\frac{L\eta_t^{\,2}}{2}\|\delta_t\|^2 .
\end{aligned}
\tag{9}
\]

\item[Step 6.] Take conditional expectation $\E[\cdot\mid w_t]$.
Using Assumption~\ref{as:unbiased}, $\E[\langle\nabla f(w_t),\delta_t\rangle\mid w_t]=0$.
Using Assumption~\ref{as:variance}, $\E[\|\delta_t\|^2\mid w_t]\le\sigma^2$.
Thus
\[
\E\!\left[f(w_{t+1})\mid w_t\right]
\le f(w_t)-\Bigl(\eta_t-\frac{L\eta_t^{\,2}}{2}\Bigr)\|\nabla f(w_t)\|^2
+\frac{L\eta_t^{\,2}}{2}\sigma^2 .
\tag{10}
\]

\item[Step 7.] Since $\eta_t<1/L$ (Assumption~\ref{as:lr}), we have
$\eta_t-\frac{L\eta_t^{\,2}}{2}\ge\frac{\eta_t}{2}$.  Hence
\[
\E\!\left[f(w_{t+1})\mid w_t\right]
\le f(w_t)-\frac{\eta_t}{2}\|\nabla f(w_t)\|^2
+\frac{L\eta_t^{\,2}}{2}\sigma^2 .
\tag{11}
\]

\item[Step 8.] Take full expectation and sum \eqref{11} over $t=0,\dots,T-1$:
\[
\begin{aligned}
\E[f(w_T)] &\le f(w_0)-\frac12\sum_{t=0}^{T-1}\eta_t\E\!\big[\|\nabla f(w_t)\|^2\big]
+\frac{L\sigma^2}{2}\sum_{t=0}^{T-1}\eta_t^{\,2}.
\end{aligned}
\tag{12}
\]

\item[Step 9.] Rearrange \eqref{12} and use $f(w_T)\ge f^\star$:
\[
\frac12\sum_{t=0}^{T-1}\eta_t\E\!\big[\|\nabla f(w_t)\|^2\big]
\le f(w_0)-f^\star+\frac{L\sigma^2}{2}\sum_{t=0}^{T-1}\eta_t^{\,2}.
\tag{13}
\]

\item[Step 10.] Divide both sides of \eqref{13} by $\sum_{t=0}^{T-1}\eta_t$
and multiply by $2$:
\[
\frac{1}{T}\sum_{t=0}^{T-1}\E\!\big[\|\nabla f(w_t)\|^2\big]
\le
\frac{2\big(f(w_0)-f^\star\big)}{\displaystyle\frac{1}{T}\sum_{t=0}^{T-1}\eta_t}
\;+\;
\frac{L\sigma^2}{\displaystyle\frac{1}{T}\sum_{t=0}^{T-1}\eta_t}
\;\frac{1}{T}\sum_{t=0}^{T-1}\eta_t^{\,2}.
\tag{14}
\]
Recognising the definitions of $\bar\eta_T$ yields exactly the bound \eqref{3}.
\end{enumerate}
All steps respect the stated assumptions, completing the proof. $\square$

\section*{Rate Derivation (With Constants)}
From the cyclic schedule \eqref{1} we have
\[
\eta_{\min}\le\eta_t\le\eta_{\max}\quad\forall t,
\]
hence
\[
\bar\eta_T\ge\eta_{\min},\qquad
\frac{1}{T}\sum_{t=0}^{T-1}\eta_t^{\,2}\le\eta_{\max}^{2}.
\]
Plugging these crude bounds into \eqref{3} yields the explicit rate \eqref{4}:
\[
\frac{1}{T}\sum_{t=0}^{T-1}\E\!\big[\|\nabla f(w_t)\|^2\big]
\le
\frac{2\big(f(w_0)-f^\star\big)}{\eta_{\min}T}
+\frac{L\sigma^2\eta_{\max}^{2}}{\eta_{\min}T}.
\]
Thus the average gradient norm decays as $\mathcal{O}(1/T)$.

\section*{Special Cases}
\begin{itemize}
\item \textbf{Constant step size.} If $\eta_{\min}=\eta_{\max}=\eta$, the bound collapses to the classical SGD result.
\item \textbf{Vanishing variance.} When $\sigma=0$ (e.g., deterministic gradients), the second term disappears and we obtain the deterministic $O(1/T)$ rate.
\item \textbf{Large period $M$.} As $M\to\infty$, the schedule becomes effectively constant, and the bound remains unchanged, confirming that the theory is agnostic to the period length.
\end{itemize}

\end{document}